\documentclass[11pt,a4paper]{article}
\usepackage[utf8]{inputenc}
\usepackage{amsmath,amssymb}
\usepackage{booktabs}
\usepackage{hyperref}

\title{Sistemas de Visión Artificial y Calibración Ojo-en-Mano para Aplicaciones Pick-and-Place}
\author{Dra. Sofía Alarcón \and Ing. Javier Mendoza \\
\textit{Laboratorio de Percepción y Visión Robótica, Centro Nova}}
\date{Junio de 2026}

\begin{document}

\maketitle

\begin{abstract}
El empaquetado y clasificación a alta velocidad en líneas industriales automatizadas requiere la sincronización milimétrica entre sensores ópticos de visión artificial y brazos robóticos SCARA o Delta (paralelos). Este artículo aborda la calibración fotogramétrica de cámaras industriales usando el modelo estenopeico (pinhole), la resolución de la ecuación matricial clásica de calibración ojo-en-mano (Hand-Eye Calibration) $AX = XB$, y la corrección de distorsión radial no lineal. Se complementa con algoritmos de seguimiento de cintas transportadoras (Conveyor Tracking) mediante encoders en cuadratura para trayectorias dinámicas de agarre.
\end{abstract}

\section{Modelo Geométrico de Cámara Estenopeica}
La proyección de un punto tridimensional en el espacio de trabajo $P_w = [X_w, Y_w, Z_w, 1]^T$ sobre el plano de píxeles del sensor CMOS $p = [u, v, 1]^T$ se describe mediante la matriz de cámara proyectiva:
\begin{equation}
s \begin{bmatrix} u \\ v \\ 1 \end{bmatrix} = K \cdot [R_{c}^w \mid t_{c}^w] \cdot \begin{bmatrix} X_w \\ Y_w \\ Z_w \\ 1 \end{bmatrix}
\end{equation}
donde $s$ es un factor de escala proyectivo, y $K$ es la matriz intrínseca de calibración:
\begin{equation}
K = \begin{bmatrix}
f_x & \gamma & c_x \\
0 & f_y & c_y \\
0 & 0 & 1
\end{bmatrix}
\end{equation}
con $f_x, f_y$ como distancias focales en píxeles, $(c_x, c_y)$ las coordenadas del punto principal óptico y $\gamma$ el factor de sesgo angular (normalmente $\gamma = 0$).

\subsection{Compensación de Distorsión de Lente}
Las lentes ópticas de gran angular inducen distorsión radial y tangencial. Las coordenadas distorsionadas $(x_d, y_d)$ se corrigen usando los coeficientes de Brown-Conrady:
\begin{align}
x_d &= x (1 + k_1 r^2 + k_2 r^4 + k_3 r^6) + 2 p_1 x y + p_2 (r^2 + 2x^2) \\
y_d &= y (1 + k_1 r^2 + k_2 r^4 + k_3 r^6) + p_1 (r^2 + 2y^2) + 2 p_2 x y
\end{align}
donde $r^2 = x^2 + y^2$, $(k_1, k_2, k_3)$ son los coeficientes radiales y $(p_1, p_2)$ los tangenciales.

\section{Calibración Ojo-en-Mano: Ecuación $AX = XB$}
En una configuración montada en el efector final (\textit{eye-in-hand}), la cámara se desplaza solidaria al extremo del robot. Para relacionar las coordenadas de los objetos observados con la base del robot, se debe determinar la transformación rígida desconocida $X = T_{cam}^{gripper}$.

Al mover el robot entre dos posturas distintas $i$ y $j$, se forma el lazo cinemático cerrado:
\begin{equation}
A \cdot X = X \cdot B
\end{equation}
donde:
\begin{itemize}
    \item $A = (T_{gripper_j}^{base})^{-1} \cdot T_{gripper_i}^{base}$ representa el desplazamiento del robot reportado por los encoders.
    \item $B = T_{target}^{cam_j} \cdot (T_{target}^{cam_i})^{-1}$ representa el desplazamiento visual medido del patrón de calibración.
\end{itemize}

\subsection{Solución Desacoplada de Tsai-Lenz}
Separando la rotación $R_X$ de la traslación $t_X$:
\begin{align}
R_A \cdot R_X &= R_X \cdot R_B \\
(R_A - I) \cdot t_X &= R_X \cdot t_B - t_A
\end{align}
Representando las rotaciones mediante vectores de eje-ángulo (o cuaterniones unitarios), se obtiene una solución de mínimos cuadrados linealmente estable para $N \ge 3$ movimientos no paralelos.

\section{Seguimiento Dinámico de Cintas (Conveyor Tracking)}
Para recoger piezas en movimiento sin detener la línea de producción, la posición del efector final en la dirección de traslación de la cinta $x_c(t)$ se sincroniza en tiempo real:
\begin{equation}
x_c(t) = x_{cam} + v_{belt} \cdot (t - t_{trigger}) + \frac{1}{2} a_{belt} (t - t_{trigger})^2
\end{equation}
donde $v_{belt}$ se computa leyendo las pulsaciones del encoder incremental de cuadratura a través del módulo de conteo rápido del PLC industrial.

\section{Conclusiones}
La integración de visión artificial calibrada milimétricamente con algoritmos de tracking cinemático permite tasas de captura superiores a 120 piezas por minuto en robots paralelos Delta, con tolerancias espaciales inferiores a $\pm 0.25\,\text{mm}$.

\begin{thebibliography}{9}
\bibitem{tsai} R. Y. Tsai and R. K. Lenz, \textit{A new technique for fully autonomous and efficient 3D robotics hand/eye calibration}, IEEE Trans. Robotics and Automation, 1989.
\bibitem{zhang} Z. Zhang, \textit{A flexible new technique for camera calibration}, IEEE Trans. Pattern Analysis and Machine Intelligence, 2000.
\bibitem{hartley} R. Hartley and A. Zisserman, \textit{Multiple View Geometry in Computer Vision}, Cambridge University Press, 2004.
\end{thebibliography}

\end{document}
